Conventions are those of Theorem \claimref{mass.exact-join} and its proof: $Z = (\omega, \pi)$ with $\omega \in \Som$, $\pi \in \Spi$; $t_{il} = \overline{\pi^{i}}\cdot\omega^{l}$ (the plain sum of products of components, the bar being the conjugation of the chart's unit, $\Sigma$ pairing $e_{3}, e_{4}$ with $e_{1}, e_{2}$); $h = t + t^{\dagger}$; $\br{u\,v} = [u\,v] = \varepsilon_{ab}u^{a}v^{b}$; $m^{2} = 2|B_{\Lambda}|^{2} + 2\lambda C$ with $C = \tfrac12\tr h^{2} - \tfrac14(\tr h)^{2}$; and by (5) of that theorem, for $\Lambda > 0$, $m^{2} = 0$ if and only if $B_{\Lambda} = 0$ and $h = 0$.

(1) If $b \neq 0$ the spinors $\pi^{1}, \pi^{2}$ are a basis of $\Spi$, so there is exactly one linear map $\mu$ with $\mu\pi^{i} = \omega^{i}$. Then $t_{il} = \overline{\pi^{i}}\cdot\mu\pi^{l}$ and $\overline{t_{li}} = \overline{\overline{\pi^{l}}\cdot\mu\pi^{i}} = \overline{\pi^{i}}\cdot\mu^{\dagger}\pi^{l}$, so $h_{il} = \overline{\pi^{i}}\cdot(\mu + \mu^{\dagger})\pi^{l}$. For a $2\times 2$ matrix, $\varepsilon(\mu u, \mu v) = \det(\mu)\,\varepsilon(u, v)$, so $\Omega_{\omega} = [\mu\pi^{1}, \mu\pi^{2}] = \det(\mu)\,b$, and $B_{\Lambda} = b + \lambda\Omega_{\omega} = (1 + \lambda\det\mu)\,b$. The span $\{\mu\pi + \pi\}$ is the graph of $\mu$; by the proof of Theorem \claimref{mass.exact-join}(4) every such graph is a line $L_{x}$, with $\mu = \mu_{x}$, $(\mu_{x}\pi)^{a} = x^{aa'}\varepsilon_{b'a'}\pi^{b'}$ and $\det\mu_{x} = \det(x^{aa'}) = -\tfrac12 Q(x)$ (the correspondence $x \mapsto \mu_{x}$ is a linear bijection of the $x$-slot onto $2\times 2$ matrices). Hence $1 + \lambda\det\mu = 1 - \Lambda Q(x)/12 = \sigma(x)$, which is (4) of that theorem read through $\mu$. If $\mu + \mu^{\dagger} = 0$ then $h = 0$, so $C = 0$ and $m^{2} = 2|B_{\Lambda}|^{2} = 2\sigma(x)^{2}|b|^{2}$; $\sigma(x)$ is real there, because $\det\mu$ of an anti-Hermitian $2\times 2$ matrix is real.

(2) Suppose $m^{2} = 0$, so $h = 0$ and $B_{\Lambda} = 0$.
\emph{Case $b \neq 0$.} By (1), $h = 0$ for the basis $\pi^{1}, \pi^{2}$ of $\Spi$ means that the Hermitian form $\overline{u}\cdot(\mu + \mu^{\dagger})v$ vanishes on a basis, so $\mu + \mu^{\dagger} = 0$; and $B_{\Lambda} = \sigma(x)b = 0$ forces $\sigma(x) = 0$. Conversely these two conditions give $h = 0$ and $B_{\Lambda} = 0$, hence $m^{2} = 0$. This is (a). The hypersurface $\sigma = 0$ is where the conformal factor of $g = \sigma^{-2}\eta$ in Lemma \claimref{mass.de-sitter-scale-frozen-point}(1) blows up: the conformal boundary of the de Sitter region $\sigma \neq 0$.
\emph{Case $b = 0$.} Then $\pi^{1}, \pi^{2}$ are dependent and the span contains a nonzero $Z_{0} = (\omega_{0}, 0)$. If every element of the span has $\pi = 0$, the span is $\Som$, $\omega^{1}, \omega^{2}$ are independent (since $Z^{1}\wedge Z^{2} \ne 0$), $h = 0$ (the form pairs $\Som$ only with $\Spi$), and $m^{2} = 2|\lambda\Omega_{\omega}|^{2} = 2\lambda^{2}|\Omega_{\omega}|^{2} > 0$: not massless. Otherwise take a basis $Z_{0}$, $Z_{1} = (\omega_{1}, \pi_{1})$ with $\pi_{1} \ne 0$; $B_{\Lambda}$ and the brackets change by the determinant of the change of basis, so their vanishing does not depend on the basis. $B_{\Lambda} = \br{0\,\pi_{1}} + \lambda[\omega_{0}\omega_{1}] = \lambda[\omega_{0}\omega_{1}]$, which vanishes iff $\omega_{1} = c\,\omega_{0}$; replacing $Z_{1}$ by $Z_{1} - cZ_{0} = (0, \pi_{1})$, the Gram matrix of the basis $(\omega_{0}, 0), (0, \pi_{1})$ has diagonal entries $0$ and off-diagonal entry $\overline{\omega_{0}}\cdot\pi_{1}$ (conjugated), so $h = 0$ iff $\overline{\omega_{0}}\cdot\pi_{1} = 0$. This is (b); conversely (b) gives $B_{\Lambda} = 0$ and $h = 0$, hence $m^{2} = 0$. For fixed $\omega_{0}$, $\overline{\omega_{0}}\cdot\pi = 0$ fixes $\pi$ up to scale, so (b) is parametrised by the line of $\omega_{0}$, a two-sphere.

\emph{Both types lie on one boundary.} By Theorem \claimref{mass.exact-join}(1), $B_{\Lambda} = \ip{I_{\Lambda}}{X}$ with $X = Z^{1}\wedge Z^{2}$, and $h = 0$ says the plane of $X$ is totally null. So for $\Lambda > 0$ the massless composites with $X \neq 0$ are exactly the totally null planes with $\ip{I_{\Lambda}}{X} = 0$. Inside the chart ($b = \ip{n_{\infty}}{X} \neq 0$) these are $X = b\,X(x)$ with $\ip{I_{\Lambda}}{X(x)} = \sigma(x) = 0$ by Lemma \claimref{mass.de-sitter-scale-frozen-point}(2): type (a). Outside it, $\ip{n_{\infty}}{X} = 0$, and then $\ip{I_{\Lambda}}{X} = \lambda\ip{n_{0}}{X} = \lambda\Omega_{\omega}$ (the computation of (1) in that theorem) must vanish too: $X \perp n_{\infty}, n_{0}$, which is type (b). So type (b) is the two-sphere of the boundary $\ip{I_{\Lambda}}{X} = 0$ at the chart's infinity. That totally null planes are the real points of compactified Minkowski space is classical twistor geometry (Penrose; cited from memory, not used beyond the name). On the whole locus $h = 0$, so the conformal Casimir $C_{2}(\Lambda^{2}\CS) = \tr h^{2} - \tfrac14(\tr h)^{2}$ vanishes and the squared grip equals the de Sitter mass there; the quantum massless locus is not addressed.

(3) On (a), $\sigma(x) = 1 + \lambda\det\mu = 0$ gives $\det\mu = -1/\lambda$ and, by (1), $\Omega_{\omega} = -b/\lambda = -(6/\Lambda)b$. Proposition \claimref{families.source-term-linear-tilt-mass-functional}(1) gives the source term $\tfrac{2\Lambda}{3}\mathrm{Re}(\overline{b}\,\Omega_{\omega})$ and its derivative in $b$, $\tfrac{2\Lambda}{3}\Omega_{\omega}$; substituting, they equal $-4|b|^{2}$ and $-4b$. (As a consistency check, the three terms of that proposition's expansion are $2|b|^{2}$, $-4|b|^{2}$ and $\tfrac{\Lambda^{2}}{18}\cdot\tfrac{36}{\Lambda^{2}}|b|^{2} = 2|b|^{2}$, summing to $2|B_{\Lambda}|^{2} = 0$.) On (b), $b = \br{0\,\pi_{1}} = 0$ and $\Omega_{\omega} = [\omega_{0}\,0] = 0$ in the adapted basis, hence in every basis. Since on (a) $b \neq 0$ and $\Omega_{\omega} \neq 0$, and on (b) both vanish, on the massless locus $b = 0$ iff $\Omega_{\omega} = 0$; at a massless composite with $b = 0$ the derivative of the source at $b = 0$, $\tfrac{2\Lambda}{3}\Omega_{\omega}$, is zero. A composite that is massless and carries a nonzero source already has $b \neq 0$ and lies on component (a). What is shown is therefore that the tilt at $b = 0$ vanishes on the massless locus wherever $b = 0$; on component (a) the source term is nonzero.

(4) By \claimref{synthesis.many-twistor-join}(v), for $\Lambda > 0$ an $N$-constituent composite is massless iff every $B^{il}_{\Lambda} = 0$ and $h = 0$. Then the span $V$ of the $Z^{i}$ is totally null for the Hermitian form of $\CS$, which has signature $(2,2)$, so $\dim V \le 2$. If $\dim V = 1$ every $Z^{i}$ is a multiple of one $Z$, and $h = 0$ says $Z$ is null; all brackets vanish. If $\dim V = 2$, choose two of the $Z^{i}$ spanning $V$; their pair is massless, so $V$ is of type (a) or (b), and every other $B^{il}_{\Lambda}$ is a multiple of theirs by bilinearity. Conversely all such configurations have $h = 0$ and every $B^{il}_{\Lambda} = 0$.

The check computes $m^{2}$ from the definition of the grip (the chart matrices of the moment map paired by an exact inverse of the ambient Gram matrix), independently of the join formula, and confirms: the identities of (1) on non-Hermitian $\mu$; $\det\mu_{x} = -\tfrac12 Q(x)$ and $B_{\Lambda} = \sigma(x)b$ for chart points; $h = 0$ and $m^{2} = 2\sigma^{2}|b|^{2}$ for anti-Hermitian $\mu$, zero exactly at $\det\mu = -1/\lambda$ and nonzero after moving $\det\mu$; the values of (3); a massless three-constituent composite in a plane of type (a); the type (b) planes massless and a tilted one massive; and $m^{2} = 2\lambda^{2}|\Omega_{\omega}|^{2}$ on the plane $\pi = 0$.
